% =====================================================================
\chapter{Evaluación: entropía cruzada y perplejidad}
\label{tema:5}
% =====================================================================

Un buen modelo asigna probabilidad alta al texto real. Este tema define la pérdida que se minimiza al entrenar todos los modelos de lenguaje, la entropía cruzada, y la medida de evaluación derivada de ella, la perplejidad. Se sigue~\cite[§3.2, §3.3 y §3.7]{jm3}; los resultados de teoría de la información proceden de~\cite[cap.~2]{ct}.

\section{Pérdida y perplejidad}

Sea un texto de evaluación con $N$ predicciones y sea $p_t=P(x_t\mid x_{1:t-1})\in(0,1]$ la probabilidad que el modelo asigna al token que realmente aparece en la posición $t$.

\begin{definicion}[Pérdida y perplejidad]\label{def:ppl}
La \emph{pérdida} (log-verosimilitud negativa media) y la \emph{perplejidad} del modelo sobre el texto son
\[
\mathcal{J}=-\frac{1}{N}\sum_{t=1}^{N}\ln p_t,\qquad \PPL=e^{\mathcal{J}}.
\]
Con logaritmo natural, $\mathcal{J}$ se mide en \emph{nats}; con $\log_2$, en \emph{bits}, y $\mathcal{J}_{\text{bits}}=\mathcal{J}_{\text{nats}}/\ln 2$.
\end{definicion}

\begin{proposicion}\label{prop:ppl}
\begin{enumerate}
  \item $\PPL=\bigl(\prod_{t=1}^{N}1/p_t\bigr)^{1/N}$, la media geométrica de los inversos de las probabilidades.
  \item $\PPL\ge 1$, con igualdad si y solo si $p_t=1$ para todo $t$.
  \item Un modelo uniforme sobre $K$ tokens tiene $\PPL=K$ sobre cualquier texto.
\end{enumerate}
\end{proposicion}

La perplejidad admite una interpretación como \emph{factor de ramificación medio}: un modelo con perplejidad $k$ está, en promedio, tan indeciso como si eligiera uniformemente entre $k$ tokens~\cite[§3.3.1]{jm3}.

\addtocounter{definicion}{1}%

\section{Entropía y entropía cruzada}

\begin{definicion}[Entropía y entropía cruzada]
Sean $p$ y $q$ distribuciones sobre un conjunto finito $\mathcal{X}$, con $q(x)>0$ siempre que $p(x)>0$. La \emph{entropía} de $p$ y la \emph{entropía cruzada} de $q$ respecto de $p$ son
\[
H(p)=-\sum_{x\in\mathcal{X}}p(x)\ln p(x),\qquad
H(p,q)=-\sum_{x\in\mathcal{X}}p(x)\ln q(x),
\]
con el convenio $0\ln 0=0$.
\end{definicion}

Si los tokens del texto se distribuyen según $p$ y el modelo es $q$, la pérdida $\mathcal{J}$ es un estimador de $H(p,q)$: es la media empírica de $-\ln q(X)$ con $X\sim p$.

\begin{teorema}[Desigualdad de Gibbs]\label{teo:gibbs}
$H(p,q)\ge H(p)$, con igualdad si y solo si $q=p$.
\end{teorema}

La diferencia $H(p,q)-H(p)=\sum_x p(x)\ln\bigl(p(x)/q(x)\bigr)$ es la \emph{divergencia de Kullback–Leibler} $D(p\,\|\,q)$, que se estudia en el módulo~6.

\begin{corolario}\label{cor:aprender}
Entre todos los modelos $q$, la entropía cruzada esperada $H(p,q)$ se minimiza exactamente en $q=p$, y su valor mínimo es $H(p)$.
\end{corolario}

El \cref{cor:aprender} justifica el criterio de entrenamiento de los modelos de lenguaje: minimizar la entropía cruzada media sobre un corpus empuja al modelo a reproducir la distribución real del token siguiente. Ningún modelo puede bajar de $H(p)$, la incertidumbre propia del lenguaje. La relación entre entropía y predicción se remonta a Shannon~\cite{shannon}.

\section{Entrenamiento y evaluación}

\begin{observacion}
La perplejidad debe medirse sobre un conjunto de evaluación disjunto del de entrenamiento. Evaluar sobre los datos de entrenamiento da una estimación optimista: el EMV del \cref{teo:emv} maximiza precisamente la verosimilitud de esos datos~\cite[§3.2]{jm3}. Además, la perplejidad solo es comparable entre modelos que usan el mismo vocabulario, porque depende de la tokenización.
\end{observacion}

